\subsection{Duomenų padalijimas}
\label{subsec:cnn_duomenu_padalijimas}

Mokymui, validavimui ir testavimui naudojamos atskiros aibės be persidengimų. Abi užduotys dalijamos su fiksuotu \textit{seed} $= 42$.

\textbf{Vaizdai.} Iš \texttt{train} aplanko stratifikuotai atskiriama $20\,\%$ vaizdų į validavimo aibę, testavimo aibė lieka nepakitusi. Po pošalio (\textit{subset}) mechanizmo gauti aibių dydžiai pateikti \ref{tab:image_split} lentelėje.

\begin{table}[H]
    \centering
    \caption{Vaizdų aibės padalijimas po validavimo atskyrimo.}
    \begin{tabular}{|l|c|c|c|}
        \hline
        Aibė & Pradinė & \texttt{subset\_ratio}~$=0{,}2$ & \texttt{subset\_ratio}~$=0{,}5$ \\
        \hline
        Mokymo & $3\,786$ & $757$ & $1\,893$ \\ \hline
        Validavimo & $947$ & $189$ & $473$ \\ \hline
        Testavimo & $1\,184$ & $236$ & $592$ \\ \hline
        \textbf{Viso} & $\mathbf{5\,917}$ & $\mathbf{1\,182}$ & $\mathbf{2\,958}$ \\ \hline
    \end{tabular}
    \label{tab:image_split}
\end{table}


\textbf{Laiko eilutės.} Atrinktoji $1\,500$ įrašų aibė skaidoma dviem stratifikuotais žingsniais (per \texttt{sklearn.model\_selection.train\_test\_split} su \texttt{stratify=y}): pirmas – $80\,\%$ mokymo + validavimo, $20\,\%$ testavimo; antras – mokymo + validavimo aibė toliau skaidoma į $80\,\%$ mokymo ir $20\,\%$ validavimo. Bendras santykis – $64{:}16{:}20$. Aibių dydžiai pateikti \ref{tab:keystroke_split} lentelėje.

\begin{table}[H]
    \centering
    \caption{Laiko eilučių aibės padalijimas (santykis $64{:}16{:}20$).}
    \begin{tabular}{|l|c|c|}
        \hline
        Aibė & Įrašų skaičius & Dalis \\
        \hline
        Mokymo & $960$ & $64\,\%$ \\ \hline
        Validavimo & $240$ & $16\,\%$ \\ \hline
        Testavimo & $300$ & $20\,\%$ \\ \hline
        \textbf{Viso} & $\mathbf{1\,500}$ & $\mathbf{100\,\%}$ \\ \hline
    \end{tabular}
    \label{tab:keystroke_split}
\end{table}


Kiekviena aibė turi savo paskirtį:
\begin{itemize}
    \item Mokymo aibė – naudojama tinklo svoriams atnaujinti gradientinio nusileidimo metu.
    \item Validavimo aibė – naudojama modelio kokybei stebėti kiekvienos epochos pabaigoje, hiperparametrams parinkti ir persimokymui (angl.~\textit{Overfitting}) aptikti.
    \item Testavimo aibė – naudojama tik galutiniam geriausio modelio įvertinimui su nematytais duomenimis.
\end{itemize}
